\section{模型假设}
\subsection{问题一：湍流强度模型的构建与优化}
在两种设备完成时间和位置匹配后，假设同一配对窗口内的廓线能够描述相近的大气状态。温度与气压采用一致单位，高度使用同一垂直基准。谱宽经过可获得的仪器、波束和切变展宽校正后，作为与湍流相关的代理量；当校正项缺失时，在指标定义和误差讨论中保留其影响。

模型 a 的输出作为题目指定的标定参照。对于训练范围内的雷达变量，假设其与参照指标之间存在可估计的映射关系。该关系的适用范围通过独立时间块检验，超出训练高度或观测范围的输入单独标记。
\subsection{问题二：三维湍流场的多源数据融合}
在选定分析时间窗内，假设三维场具有有限的空间相关尺度，且水平和垂直方向允许采用不同参数。设备误差以协方差描述；只有在误差相关性可以忽略时才使用对角近似。缺少观测的区域由背景项和平滑项共同约束，同时输出覆盖信息及误差估计。

地面站的约束限制在其代表高度附近。对于雷达未覆盖的低层或遮挡区，使用相应缺测掩膜，避免把没有观测解释为风险为零。平滑约束表达空间连续性，质量和能量守恒需要由独立动力方程及边界条件检验。
\subsection{问题三：低空飞行的航路规划与优化}
假设起点、终点和允许飞行区域来自题目规定，飞行速度或到达时间规则在寻路前确定。风险代价在各边上非负，图中的每条边均通过沿线空域检查。预测场的时间插值只读取该次预报起点已经生成的场，不接入未来观测。

题面要求的主模型以湍流条件作为优化目标。航程、转弯和爬升可作为可行性约束或扩展方案，采用扩展方案时另列目标与对照结果。表\ref{tab:assumptions}列出了假设对应的检验方式。
\begin{table}[H]\centering\auxtablecaption{模型假设的检查与影响分析}\label{tab:assumptions}
\begin{tabularx}{\linewidth}{YYY}\toprule
假设 & 检查内容 & 偏离时的处理\\\midrule
配对窗口内状态相近 & 改变时间容差后的廓线差异 & 缩短窗口或扩大误差\\
空间相关尺度有限 & 方向半方差和留站误差 & 分区设定尺度\\
误差可用协方差表达 & 同设备残差相关性 & 保留非对角项或分组重采样\\
风险可沿路径累计 & 网格加密后的积分变化 & 细分边并复算\\
预测信息满足时间边界 & 起报时刻与输入时间检查 & 删除越界输入并重新训练\\\bottomrule
\end{tabularx}\end{table}
\auxfigure{assumptions}{假设、参数扰动与验证结果的对应关系}{fig:assumptions}
